\documentclass{kdiartcl}

\usepackage{amsmath}

\usepackage{tabularx}
\usepackage{array}
\usepackage{ragged2e}
\usepackage{csquotes}

\newcolumntype{Y}{>{\RaggedRight\arraybackslash}X}

\usetikzlibrary{arrows.meta,decorations.pathreplacing}

\university{HTWG Konstanz}
\institute{Fachbereich Angewandte Informatik}
\lecture{Rechnernetze}
\date{\today}

\author{Paul Krüssel}
\title{Lösung Praxisaufgabe 3}
%\subtitle{Abgabeende: 14. April 2026}

\singlespacing

\begin{document}
\maketitle

\setcounter{section}{1}

\section{Monitoring von Sockets}
\subsection{Wieviele Sockets sind insgesamt geöffnet?}
Ich habe das Tool CurrPorts auf Windows genutzt. Es zeigt 167 Current Ports und 10 Remote Connections an.
\subsection{Wie unterscheiden sich die Einträge von TCCP und UDP Sockets?}
Die TCP Ports haben verschiedene States, während UDP keine States hat.
\subsection{Was bedeuten die Einträge in der Spalte "State" bei TCP Sockets?}
Zum Zeitpunkt der Analyse habe ich diese States gesehen:
\begin{itemize}
    \item LISTEN: Der Port ist geöffnet und wartet auf einen Verbindungsversuch von einem anderen TCP Port.

    \item ESTABLISHED: Der Port hat sich mit einem anderen Port verbunden und die Verbindung wird aktiv hochgehalten.
 
    \item CLOSE WAIT: Wartete auf eine Beendigungsanfrage durch einen lokalen Prozess, bzw. den lokalen User.

    \item TIME WAIT: Der Port wurde geschlossen, der Port wartet noch auf die letzten Pakete.
\end{itemize}
Es gibt aber zusätzlich noch diese States:
\begin{itemize}
    
    \item SYN-SENT: Wartet auf eine passende Verbindungsanfrage, nachdem zuvor selbst eine Verbindungsanfrage gesendet wurde.
    
    \item SYN-RECEIVED: Wartet auf eine bestätigende Verbindungsanfrage, nachdem sowohl eine Verbindungsanfrage empfangen als auch gesendet wurde.
    
    \item FIN-WAIT-1: Wartet auf eine Verbindungsbeendigungsanfrage von der entfernten TCP-Gegenstelle oder auf eine Bestätigung der zuvor gesendeten Verbindungsbeendigungsanfrage.
    
    \item FIN-WAIT-2: Wartet auf eine Verbindungsbeendigungsanfrage von der entfernten TCP-Gegenstelle.
    
    \item CLOSING: Wartet auf eine Bestätigung der Verbindungsbeendigungsanfrage von der entfernten TCP-Gegenstelle.
    
    \item LAST-ACK: Wartet auf eine Bestätigung der zuvor an die entfernte TCP-Gegenstelle gesendeten Verbindungsbeendigungsanfrage, die auch eine Bestätigung ihrer Verbindungsbeendigungsanfrage enthält.
    
    \item CLOSED: Es besteht kein Verbindungszustand.
\end{itemize}

Es sind 40 Ports auf LISTENING offen. Zudem sind 47 UDP Ports offen.

\subsection{Wie viele Sockets (ESTABLISHED) werden neu geöffnet, wenn Sie die Messung nach einer Minute
erneut durchführen bzw. die Ergebnisse aktualisieren?}
Erste Messung: 150 Ports offen, davon sind 56 auf dem State ESTABLISHED \\
Messung eine Minute später: Es sind 2 neue Sockets auf dem State ESTABLISHED, allerdings hat sich ein Port den State auf CLOSE WAIT geändert. Es sind nun isngesamt 153 Ports offen, davon sind 57 auf dem State ESTABLISHED

\subsection{Sehen Sie zahlreiche Sockets mit IP-Adresse 127.0.0.1? Finden Sie heraus, wofür diese IP Adresse
benutzt wird und blenden Sie alle Sockets mit dieser Adresse aus.}
Es sind 37 Ports offen, die $127.0.0.1$ entweder als Lokale oder als Remote Adresse nutzen (oder beides). $127.0.0.1$ ist eine lokale Loopback Adresse. Sockets können sich also mit anderen Sockets, welche auch auf der lokalen Maschine laufen, über diese IP-Adresse verbinden.\\
Der Filter 
\begin{itemize}
    \item exclude:both:tcpudp:127.0.0.1
\end{itemize}
blendet alle TCP und UDP Sockets aus, welche diese IP-Adresse nutzen.

\subsection{Bestimmen sie anhand der Portnummer und der Portliste für einige interessante/unbekannte
Prozesse, mit welchem Protokoll diese kommunizieren. }
\begin{itemize}
    \item Ein Systemprozess hat 5 Ports offen. Zwei davon akzeptieren Verbindungen über den lokalen Port $445$, welcher für das MS-Active Directory, also Windows Ordner, welche der Nutzer an das lokale Netzwerk freigegeben hat, da ist. Da diese Sockets aber auf LISTENING stehen, greift aktuell kein lokaler Nutzer auf einen freigegebenen Ordner zu.
    \item Spotify.exe hat mehrere Ports geöffnet, welche als Remote Port $443$ haben, was bedeutet, dass diese per HTTPS-Protokoll mit dem Spotify Server verbunden sind, um z.B. Informationen über Lieder aufzurufen. Diese Verbindung ist aktiv, da der State auf ESTABLISHED ist.
    \item Auch Discord.exe hat eine aktive HTTPS-Verbindung offen mit dem State auf ESTABLISHED. Insgesamt sind gerade 43 Sockets mit einer HTTPS-Verbindung offen, wobei nicht alle ESTABLISHED sind.
    \item Spotify.exe hat zudem noch einen UDP-Socket mit dem Port $1900$ auf der IP-Adresse $0.0.0.0$ offen. Dieser Port ist für das SSDP, also das Simple Service Discovery Protocol, welches ermöglicht, andere Geräte im Netzwerk zu finden, bzw. andere Geräte im Netzwerk sehen, dass dieser Prozess Netzwerkgeräte annehmen möchte. Bei Spotify wird das wahrscheinlich dem Zweck dienen, sich mit Netzwerk-Lautsprechern zu verbinden.
\end{itemize}

\section{Details über die Kommunikationspartner ihres PCs}
\subsection{IPNetInfo und Wireshark}
Das Programm IPNetInfo kontaktiert einen Arin WHOIS Server auf der IP $199.5.26.46$. Dazu wird das WHOIS Protokoll genutzt.
\subsection{Im welchem Netz befindet sich der Web-Server, der in der ersten WireShark-Aufgabe aufgerufenen
Web-Seite?}
Die Webseite aus Aufgabe 1 ist diese: http://gaia.cs.umass.edu/wireshark-labs/INTRO-wireshark-file1.html \\
Der Webserver liegt im UMASS-NET Netzwerk, welches zur university of Massachusetts gehört.
\subsection{Welche Informationen finden Sie über die HTWG?}
Ich bin in Firefox auf die Website der Hochschule gegangen und habe währenddessen CurrPorts gelaufen. Ich habe diese Infos in IPNetInfo gefunden:
\begin{itemize}
    \item IP-Adress: 2001:7c0:5f0:f020::20:31
    \item Status: Succeed
    \item Country: Germany
    \item Network Name: HTWG-Konstanz
    \item Owner Name: UNBEKANNT
    \item CIDR: 2001:7c0:5f0::/48
    \item Allocated: Yes
    \item Contact Name: Michael Steuert
    \item Address: Fachhochschule Konstanz, Hochschule fuer Technik, Wirtschaft und Gestaltung, Brauneggerstrasse 55, D-78462 Konstanz, Germany
    \item Email: steuert@htwg-konstanz.de
    \item Abuse Email: abuse@htwg-konstanz.de
    \item Abuse Contact: abuse@belwue.de
    \item Phone: +49 7531 206 434
    \item Fax: +49 7531 206 153
    \item Whois Source: RIPE NCC
\end{itemize}		

\section{Sockets beim Laden einer Web-Seite}
\subsection{Anzahl geöffneter Sockets Spiegel}
\subsubsection{Mit Ad-Blocker und Privacy Blocker}
Als AdBlocker nutze ich die Firefox Extension uBlock Origin, als Privacy Blocker nutze ich die Extension Privacy Badger
Vor dem Öffnen von Spiegel.de sind bereits 17 Sockets durch Firefox geöffnet, wahrscheinlich um die Artikel auf der Startseite von Firefox anzuzeigen.
Mit den Blockern öffnet Firefox durch das Aufrufen von spiegel.de insgesamt 4 TCP Ports.
\subsubsection{Ohne Ad- und Privacy Blocker}
Ich habe spiegel.de nun in einem privaten Fenster geöffnet, wodurch die Blocker aus sind. Nun werden 24 neue Ports durch die Verbindung mit spiegel.de geöffnet.
\subsection{Was ist die maximale Anzahl von Sockets pro Remote-IP-Adresse?}
Die IP-Adresse mit den meisten Sockets ist $136.243.25.9$ mit 4 offenen Sockets.
\subsection{Welche Remote-Ports werden verwendet?}
\begin{itemize}
    \item 136.243.25.9: 5 Sockets
    \item 34.107.221.82: 1 Socket
    \item 34.107.243.93: 2 Sockets
    \item 34.120.208.123: 1 Socket
    \item 63.140.62.120: 1 Socket
    \item 63.140.62.236: 1 Socket
    \item 128.65.210.181: 1 Socket
    \item 2600:1901:0:38d7:: : 2 Sockets
    \item 2600:9000:20ae:6400:a11c:1180:93a1 : 1 Socket
    \item 2600:9000:20ba:4000:f2aa:1fc0:93a1 : 1 Socket
    \item 2600:9000:2165:4600:337c:93cc:93a1 : 1 Socket
    \item 2600:9000:2165:fa00:1:9777:c740:21 : 1 Socket
    \item 2a02:26f0:3500:591::1e80 : 1 Socket
    \item 2a02:26f0:480:33::212:40c8 : 1 Socket
    \item 2a02:26f0:480:33::212:40df : 2 Sockets
    \item 2a04:4e42:6f::347 : 3 Sockets
    \item 2a04:4e42:8d::347 : 5 Sockets
    \item 2a05:d018:1b59:1864:84a6:445d:bd03:b26d : 1 Socket
    \item 2003:8:10:2:0:a10:985f:325a : 1 Socket
    \item 2400:52e0:1e00:2::1329:1 : 1 Socket
\end{itemize}

\subsection{Analyse IPNetInfo}
Insgesamt gab es 20 individuelle IP-Adressen. Davon konnten bei 3 kein Contact ermittelt werden. Aus den restlichen 17 Adressen lassen sich aber insgesamt 10 unterschiedliche Firmen zuordnen. Diese sind:
\begin{itemize}
    \item Adobe Inc.
    \item Amazon Data Services Ireland Techical Role Account
    \item Amazon.com, Inc.
    \item BUNNYWAY Administrator
    \item Fastly RIR Administrator
    \item Google LLC
    \item Hetzner Online GmbH
    \item Link11 GmbH Hostmaster
    \item Akamai Technologies
\end{itemize}

\subsection{Andere Populäre Website}
Als Gegenanalyse habe ich nun zeit.de genommen, welche 22 Sockets geöffnet hat. Darin waren 16 individuelle IP-Adressen. Dies sind die Firmen:
\begin{itemize}
    \item Adobe Inc.
    \item Amazon Data Services Ireland Techical Role Account
    \item Cloudflare, Inc.
    \item Fastly RIR Administrator
    \item Fastly, Inc.
    \item Google LLC
    \item Hetzner Online GmbH
    \item Akamai Technologies
    \item Akamai International B.V.
    \item RIPE Network Coordination Center
\end{itemize}

Diese Firmen tauchen nur bei spiegel.de auf:
\begin{itemize}
    \item Amazon.com, Inc.
    \item BUNNYWAY Administrator
    \item Link11 Hostmaster
\end{itemize}

Diese Firmen tauchen nur bei zeit.de auf:
\begin{itemize}
    \item Cloudflare, Inc.
    \item RIPE NCC
\end{itemize}

Diese Firmen tauchen bei beiden auf:
\begin{itemize}
    \item Adobe Inc.
    \item Amazon Data Services Ireland Techical Role Account
    \item Fastly RIR Administrator
    \item Fastly, Inc.
    \item Google LLC
    \item Hetzner Online GmbH
    \item Akamai Technologies
    \item Akamai International B.V.
\end{itemize}

Ich denke, dass die Firmen, die sich bei den beiden Seiten unterscheiden, für Hosting und Namespace zuständig sind, während die sich überlappenden Firmen wahrscheinlich für allgemeinere Dinge, wie z.B. Analytics oder Werbung da sind.

\end{document}

